\documentclass[12pt]{article}
\usepackage[utf8]{inputenc}
\usepackage[T1]{fontenc}
\usepackage[margin=1in]{geometry}
\usepackage{hyperref}
\usepackage{lmodern}
\hypersetup{colorlinks=true,linkcolor=blue,citecolor=blue,urlcolor=blue}

\usepackage{setspace}

\begin{document}
\begin{titlepage}
\centering
\vspace*{2cm}
{\LARGE\bfseries Election Forensics in a Polarized Democracy:\\[10pt]
Evidence from Peru's 2021 Presidential Runoff\par}
\vspace{1.8cm}
{\large Carlos C\'esar Ch\'avez Padilla\thanks{%
  Center for the Economics of Human Development, University of Chicago.
  Email: \href{mailto:cchavezp@uchicago.edu}{cchavezp@uchicago.edu}.
  \textit{Acknowledgments:}
  I thank the Oficina Nacional de Procesos Electorales for making
  mesa-level data publicly available. All errors are my own.
  \textit{Data availability:} The primary data are publicly available
  at \url{https://datosabiertos.gob.pe}. Replication code is available
  at \url{https://github.com/cesarchavezp29/peru-election-forensics-2021}.}\par}
\vspace{0.5cm}
{\normalsize September 2024\par}
\vspace{1.2cm}

\begin{minipage}{0.88\textwidth}
\small
\begin{center}\textbf{Abstract}\end{center}
\noindent
Peru's 2021 presidential runoff between Pedro Castillo and Keiko Fujimori
produced the closest result in modern Peruvian history---a margin of 44,263
votes out of 17 million cast. Fujimori refused to concede, alleged systematic
fraud, and delayed certification by six weeks. I apply multiple forensic
methods to the complete universe of polling-station records published by
Peru's electoral authority (ONPE): distributional fingerprint analysis,
last-digit uniformity tests, and ecological regression comparing first- and
second-round results. I also exploit the institutional distinction between
\emph{contested} actas---those resolved by the electoral tribunal after
challenge---and uncontested actas as a within-election test. No method
detects patterns consistent with vote manipulation. Contested and uncontested
polling stations show identical forensic signatures once geographic location
is taken into account. The data instead reveal extreme geographic polarization:
Castillo received over 80\,\% of valid votes in southern highland departments
while Fujimori dominated coastal and metropolitan areas.

\medskip
\noindent\textbf{Keywords:} election forensics, electoral fraud, Peru,
geographic polarization, distributional fingerprints, digit tests.

\medskip
\noindent\textbf{JEL codes:} D72, C12, P16.
\end{minipage}

\end{titlepage}
\end{document}
